\chapter{Konventionen und Beispiele}
Es folgen ein paar nützliche Beispiele und Konventionen.

\section{Gleichungen} \label{sec:gleichungen}
Bei umfangreichen Dokumenten sollte eine Gleichung, wie
%
\begin{equation} \label{equ:systemmatrix}
	\mathbf{A}_\mathrm{s} = \begin{bmatrix}1 & 0 \\ 0 & 2\end{bmatrix}
\end{equation}
%
nur eine Nummer bekommen, wenn sie, wie Gleichung~\eqref{equ:systemmatrix}, im Text referenziert wird. Die Nummerierung von Umgebungen kann über den * gesteuert werden.
% Absätze vor und nach Gleichungen vermeiden, um das Dokuemnt nicht unnötig auseinanderzuzerren

Gleichungen sind Teil des Satzes und bekommen deshalb auch Satzzeichen, wie beispielsweise die wahllosen Werte
%
\begin{align}
	\beta &= 5^{\circ}, \nonumber  \\
	c &= 42 \label{equ:antwort} 
\end{align}
%
und
\begin{equation*}
	\dot{\hat{x}} = 7 \, \mathrm{\frac{m}{s}}.
\end{equation*}
Wie an Gleichung~\eqref{equ:antwort} zu sehen, können mehrere Gleichungen im Block ausgerichtet und einzeln nummeriert und referenziert werden.

\section{Bilder und Tabellen}
Außer bei Fotos sollte auf Pixelgrafiken verzichtet werden. Vektorgrafiken können beispielsweise als PDF eingebunden werden, wie Abbildung~\ref{fig:rtm}, oder direkt mit TikZ erstellt werden.

\begin{figure} %[h] die Figure-Umgebung bietet eine Option zur Steuerung der Platzierung
	\centering
	\includegraphics[width=0.4\textwidth]{rtm_mit_schrift.pdf}
	\caption{Regelungstechnik und Mechatronik}
	\label{fig:rtm}
\end{figure}

Wie in Tabelle~\ref{tab:tabelle} zu sehen, bekommen Tabellen eine Überschrift.
%
\begin{table}
	\centering
	\caption{Tabellenüberschrift}	
	\begin{tabular}{r|cc}
		& A & B \\
		\hline
		C & 1 & 2 \\
		D & 3 & 4 \\
	\end{tabular}
	\label{tab:tabelle}
\end{table}
%

\section{Blockschaltbilder}
Blockschaltbilder können mit der im Hauptdokument eingebundenen TikZ-Bibliothek erstellt werden.
%
\begin{center} % wird von LaTex nicht frei am Seitenanfang/Ende platziert. Im Normalfall aber besser figure nutzen.
%	\centering
	\begin{tikzpicture}[node distance=5mm and 10mm]
	\node[terminal]	(w)	{};
	\node[sum]			(s) [right=of w]	{};
	\node[block]		(Regler) 		[right=of s] 			 	{Regler};
	\node[block]		(Aktor)	 		[right=of Regler]  	{Aktor};
	\node[block]		(Strecke)		[right=of Aktor]	 	{Strecke};
	\node[branch]		(punkt)			[right=of Strecke] 	{};
	\node[terminal]	(y)					[right=of punkt]		{};
	\node[block]		(Messglied) [below=of Aktor]		{Messglied};
	
	
	\draw[to] (w) -- (s) node[near start,above] {$w$};
	\draw[to] (s) -- (Regler) node[midway,above] {$e$}; 
	\draw[to] (Regler) -- (Aktor);
	\draw[to] (Aktor) -- (Strecke);
	\draw[to] (Strecke) -- (y) node[pos=0.8,above] {$y$};
	\draw[to] (punkt) |- (Messglied);
	\draw[to] (Messglied) -| (s) node[pos=0.95, right] {-};
\end{tikzpicture}
	\captionof{figure}{TikZ Blockschaltbild in Figure-Umgebung}
	\label{fig:bsb_beispiel}
\end{center}

Variablen, wie $w$ aus Abbildung~\ref{fig:bsb_beispiel}, kommen auch im Fließtext immer in eine Mathe-Umgebung.

\section{Plots}
Plots können direkt in \LaTeX mit dem Paket \texttt{pgfplots} sogar mit der TU-Farbpalette gesetzt werden.
In Abbildung~\ref{fig:step_pt2} wird z.B. die rtm-Farbe \textcolor{TUDa-2c}{\texttt{TUDa-2c}} genutzt.

Aus Matlab heraus können Plots z.B. mit dem Tool \texttt{matlab2tikz} für \LaTeX{} exportiert werden.
Der mitgelieferte \texttt{cleanfigure()}-Befehl räumt Plots vor dem Export auf und sorgt für eine kleine Dateigröße.
Mit \texttt{pgfplots} kann aber auch direkt in \LaTeX analytische Funktionen plotten: 
\tikz{	
	\begin{axis}[hide axis, height=18mm, width=5cm]
		\addplot[domain=-2*pi:2*pi, samples=100]{cos(deg(x))};
	\end{axis}
}.

\begin{figure}
	\centering
	% !TeX root = ../TUDaThesis.tex

% This file was created by matlab2tikz.
%
%The latest updates can be retrieved from
%  http://www.mathworks.com/matlabcentral/fileexchange/22022-matlab2tikz-matlab2tikz
%where you can also make suggestions and rate matlab2tikz.
%
\definecolor{mycolor1}{rgb}{0.00000,0.44700,0.74100}%
%
\begin{tikzpicture}

\begin{axis}[%
width=0.7\textwidth,
height=6cm,
scale only axis,
separate axis lines,
every outer x axis line/.append style={white!40!black},
every x tick label/.append style={font=\color{white!40!black}},
every x tick/.append style={white!40!black},
xmin=0,
xmax=14,
every outer y axis line/.append style={white!40!black},
every y tick label/.append style={font=\color{white!40!black}},
every y tick/.append style={white!40!black},
ymin=0,
ymax=2.5,
ylabel={Ausgang normiert},
xlabel={Zeit in Sekunden},
axis background/.style={fill=white}
]
\addplot [color=TUDa-2c, line width=1, forget plot]
  table[row sep=crcr]{%
0	0\\
0.0921034037197206	0.00411279805587128\\
0.184206807439441	0.0159482926069074\\
0.276310211159164	0.0347790289943219\\
0.368413614878884	0.0599132256263424\\
0.460517018598605	0.0906944372784881\\
0.552620422318325	0.126501104885605\\
0.644723826038048	0.166746003636801\\
0.736827229757768	0.210875600557069\\
0.828930633477489	0.258369332141982\\
0.921034037197209	0.308738812005936\\
1.01313744091693	0.361526977911547\\
1.10524084463665	0.416307186968996\\
1.19734424835637	0.472682267230876\\
1.38155105579581	0.588769773524355\\
1.93417147811414	0.942361883863663\\
2.11837828555358	1.05580915778143\\
2.2104816892733	1.11091852101091\\
2.30258509299302	1.16478018774458\\
2.39468849671275	1.21728576867211\\
2.48679190043247	1.26834190455422\\
2.57889530415219	1.31786932190199\\
2.67099870787191	1.36580191469028\\
2.76310211159163	1.41208585353384\\
2.85520551531135	1.45667872351849\\
2.94730891903107	1.49954869165797\\
3.03941232275079	1.54067370474229\\
3.13151572647051	1.5800407181526\\
3.22361913019023	1.61764495604156\\
3.31572253390996	1.65348920311623\\
3.40782593762968	1.68758312811146\\
3.4999293413494	1.71994263890565\\
3.59203274506912	1.75058926910707\\
3.68413614878884	1.77954959582612\\
3.77623955250856	1.80685468824809\\
3.86834295622828	1.83253958652982\\
3.960446359948	1.85664281046295\\
4.05254976366772	1.87920589727478\\
4.14465316738744	1.90027296787515\\
4.23675657110716	1.91989032080341\\
4.32885997482689	1.93810605308292\\
4.42096337854661	1.95496970715117\\
4.51306678226633	1.97053194300157\\
4.60517018598605	1.98484423464626\\
4.69727358970577	1.99795858998971\\
4.78937699342549	2.00992729318779\\
4.88148039714521	2.02080266855745\\
4.97358380086493	2.03063686509711\\
5.15779060830437	2.04738828496254\\
5.34199741574382	2.06058825850954\\
5.52620422318326	2.07063002123016\\
5.7104110306227	2.07788951993463\\
5.89461783806214	2.08272250421635\\
6.07882464550158	2.08546241975432\\
6.26303145294103	2.08641900115247\\
6.44723826038047	2.08587746840164\\
6.72354847153963	2.08281911140062\\
6.99985868269879	2.07774554823536\\
7.36827229757768	2.06896203894915\\
8.19720293105517	2.04625070133987\\
8.74982335337349	2.03209506210472\\
9.2103403719721	2.02202542105953\\
9.6708573905707	2.0138167145694\\
10.1313744091693	2.00747359450982\\
10.5918914277679	2.00284185840762\\
11.1445118500862	1.99919994883539\\
11.7892356761243	1.99698913404195\\
12.6181663096018	1.99626757654961\\
13.9997173653976	1.99742653315299\\
14.0918207691173	1.99753717612361\\
};
\addplot [color=black, dotted, forget plot]
  table[row sep=crcr]{%
-1	2\\
15.4	2\\
};
\end{axis}

\end{tikzpicture}%
	\caption{Sprungantwort $\mathrm{PT}_2$}
	\label{fig:step_pt2}
\end{figure}